\section{Training schedule for the M1 estimator}
\label{sec:schedule}

\Cref{sec:pseudoinj} fixes the network and the loss of \cref{eq:pseudoinj-loss}. This
section fixes \emph{when} each loss term is switched on, when the graph is allowed to
commit, how a run is monitored and stopped, and which part of the code decides what.
The reason a schedule is needed at all: the terms conflict early in training (the KL
term pulls the posterior onto an untrained prior, the sparsity term deletes edges
before the latents mean anything, the BLAE terms need a decoder that already
reconstructs), and the sink check is irreversible.

Most of the schedule is of one kind: a loss term is switched on later, by a ramp, once
something else is ready. One decision is of a different kind. The requirement that the
content graph is a DAG with a fixed causal order is not a loss term with a weight; it
is a structural constraint of the model, and the question is \emph{when to begin
enforcing it}. This is treated separately in \cref{sec:schedule:phase0}.

Every number below is a starting guess, to be fixed from the first runs. Nothing here
has been run on real TI data or on the real IN-VAE checkpoint, and the torch-dependent
code has not been run at all; only syntax and torch-free logic (for example
\texttt{config.py}) have been checked.

\subsection{Code layout and who decides what}
\label{sec:schedule:layout}
The code is split so that each file has one job. All tunable numbers live in
\texttt{config.py} (class \texttt{RunConfig}); no other file hard-codes a default.
\begin{center}
\begin{tabular}{@{}ll@{}}
\toprule
file & job\\
\midrule
\texttt{config.py} & the configuration tree \texttt{RunConfig}: \texttt{loss},
\texttt{gate}, \texttt{schedule}, \texttt{probe}\\
\texttt{head.py} & trainable head: \texttt{AutoencoderKL} plus the two linear
layers\\
\texttt{m1net.py} & head plus M1 prior, adjacency, order blend, sink-freezing and the
sink-check gate\\
\texttt{losses.py} & every loss computation; each loss owns its weight mechanics\\
\texttt{model.py} & \texttt{M1Model}: network plus losses, and the global cross-loss
decisions\\
\texttt{trainer.py} & loops, optimizers, validation, checkpoints, logging, failure
log\\
\texttt{probe.py} & gradient-scale probe (diagnostic only)\\
\texttt{main.py} & data, split, geodesic table, assembly, resume\\
\bottomrule
\end{tabular}
\end{center}
The division of authority is strict.
\begin{itemize}[leftmargin=1.5em]
\item \textbf{Each loss} knows how to compute itself, its target weight and ramp, its
      own state (started, ramp position, switched off) and whether it is
      \emph{ready} for others to wait on.
\item \textbf{The prior} (\texttt{M1Prior}) owns the structural state: the adjacency,
      the freezing buffers, the sink-check gate and the order blend of
      \cref{sec:schedule:phase0}.
\item \textbf{The model} makes only the decisions that involve more than one object
      (``do not start KL before REC has plateaued'', ``start SPARSE when KL has
      settled''), one method per rule, and says \emph{stop} when a rule requires it.
      It never computes a loss and never owns a loop.
\item \textbf{The trainer} owns every loop and every call, and executes ``stop''.
\end{itemize}

\subsection{General decisions}
\label{sec:schedule:general}
\begin{itemize}[leftmargin=1.5em]
\item \textbf{One global switch.} \texttt{schedule.enabled} is the only switch. When it
      is off, every loss uses its target weight from epoch 1, the order blend is $1$
      from epoch 1 (the lower-triangular design is imposed at once, as in the
      reference code of \cref{app:code}), the sink check runs on the fixed clock of
      \cref{app:code} (every \texttt{gate.check\_epoch} epochs), there are no caps and
      no fine-tune phase, and the run stops only at \texttt{max\_epochs}. Two fields
      are derived from it by \texttt{RunConfig.finalize()} whatever the user wrote:
      \texttt{loss.schedule} and \texttt{gate.mode} (\texttt{"stable"} or
      \texttt{"clock"}).
\item \textbf{Start big.} Begin with $n_c$ of $4$ or $5$ or more, exceeding the budget
      of \cref{sec:reqs} if necessary, and scale down only if runtime or accuracy
      demands it. $D=n_c+n_s$ with $n_s\le2$.
\item \textbf{Validation} is a held-out slice of the training locations, never a
      held-out location and never L100. It is stratified by the $(y,t,e)$ cell
      (\texttt{--val\_frac}, default $0.1$; a cell with fewer than two rows goes
      entirely to training). If synthetic images are included the split is
      \emph{not} by \texttt{src\_uuid}, so a source image and its transported copies
      can straddle the split; \texttt{main.py} prints a warning (deferred).
\item \textbf{Graph signals are logged, never pass/fail}: edge count, number of flips
      of the binarized block, margin of the raw entries to $0.5$, frozen nodes with
      the epoch of freezing, and degrees. Of these, the code currently records the
      flip history, the freeze log (epoch and node), the final adjacency with parent
      and child counts, and, at a gate cap, the five raw entries nearest $0.5$.
\item \textbf{Dimension dependence is not linear.} Candidate edges grow as
      $m(m-1)/2$, KL and BILIP grow with $D$, while REC is a mean. Only SPARSE is
      normalized (divided by the fixed constant $n_c(n_c-1)/2$); normalizing the rest
      is deferred until runs at different $n_c$ exist.
\item \textbf{Cap policy.} A rule whose effect is reversible continues past its cap
      with a warning (BILIP; KL if the head is learning; the settle wait of
      \cref{sec:schedule:phase0}, after which SPARSE starts anyway). The decisions
      that cannot be repaired stop the run: the sink-check cap, and the KL cap when
      the head is not learning. Any cap hit makes the run suspect.
\item \textbf{Timing.} A decision made at the end of epoch $e$ takes effect from epoch
      $e+1$. Weights are read on every forward pass, so a ramp started at the end of
      epoch $e$ uses its initial weight in epoch $e+1$ and gains one step at the end
      of each following epoch. The order blend of \cref{sec:schedule:phase0} follows
      the same timing.
\end{itemize}

\subsection{Phase 0: when to begin enforcing the DAG}
\label{sec:schedule:phase0}

\paragraph{The problem.}
The content graph is a strictly lower-triangular adjacency: \texttt{torch.tril} with
\texttt{diagonal=-1} is applied to the learnable matrix, so latent slot $j$ can be a
parent of slot $i$ only if $j<i$, and the causal order is the slot order. This is
inherited from the reference code of \cref{app:code}, where it is applied from epoch 1,
and there is no principled reason for that timing. Three things are wrong with it for
a real run.
\begin{itemize}[leftmargin=1.5em]
\item \textbf{The order is imposed before the slots mean anything.} At epoch 1 the head
      is random, so slot $i$ is not yet a particular factor of the data. Declaring
      ``slot 0 is a root, slot $n_c-1$ may depend on everything'' at that moment means
      that the roles are assigned by the accident of early dynamics, while the KL term
      starts to pull the posterior towards a prior that already has this asymmetric
      structure.
\item \textbf{It is a hard requirement, not a preference.} A loss weight can be ramped
      from small to large. A structural constraint cannot be made ``a little bit
      active'' by a weight; it is either in the model or not. This is why the DAG
      requirement cannot be scheduled the way SPARSE or BILIP are.
\item \textbf{The order is fixed, not searched.} The \emph{true} DAG can be expressed by a 
      lower-triangular adjacency under some assignment of latents to slots. But the training loss 
      might not select that assignment: for example, a chain and a fork over the same skeleton
      (undirected graph) fit the data equally well.
\footnote{Example with three latents.
\begin{itemize}
\item Let the true graph be the chain $Z_1\to Z_2\to Z_3$. 
If the middle node $Z_2$ is placed during learning in slot 1, the model cannot use the
edge $Z_1\to Z_2$, since it would run from a higher slot 2 to a lower slot 1. It instead
fits the fork $Z_1\leftarrow Z_2\to Z_3$. The chain and the fork have the same
skeleton and are Markov-equivalent, so with flexible mechanisms they describe the
same distributions. The likelihood is therefore the same, the loss has no reason to
object, and the model ends up with a DAG that is not the true one. It is not even
isomorphic to it: the chain has a path of length two, while the fork has a node with
two children, so no renaming of nodes turns one into the other. 
\item A collider
$Z_1\to Z_2\leftarrow Z_3$ behaves differently. It is not Markov-equivalent to the
chain or the fork, and with the $Z_2$ collider node in slot 1 a lower-triangular graph
needs an extra edge between $Z_1$ and $Z_3$ to fit the same distribution. The
sparsity penalty works against that extra edge, so some wrong orders are penalized,
but the chain-to-fork case is not.
\end{itemize}
}%
      Note that there're two potential mismatches with the true graph here:
      \begin{enumerate}
	  \item \label{item:perm} \textbf{Relabelling the slots of the true graph.} The \emph{learned} 
	     graph is the true DAG with its nodes renamed -- the two graphs are isomorphic as directed 
	     graphs. This is \emph{harmless} -- it's just a \emph{permutation} of the true DAG.
      \item \textbf{A \emph{different} DAG that fits equally well.} The \emph{learned} graph has 
         some edges reversed but the same skeleton. This is not a relabeling of the true DAG as in
         (\labelcref{item:perm}) above -- the two graphs are not isomorphic as directed graphs and 
         it is \emph{harmful}.
      \end{enumerate}       
      A wrong order shows up as a \emph{false sink}, a node frozen although it has a
      child (see \cref{sec:apply:falsesink}); it mostly does not show in the losses
      and the run ends as \texttt{finished}. Delaying the constraint does not remove
      this, but it at least lets the order be imposed on latents that carry
      information.
\end{itemize}

\paragraph{The solution: an order blend.}
The prior holds a scalar $a\in[0,1]$, the \emph{order blend}, and uses the adjacency
\begin{equation}
A_a=(1-a)\,A_{\mathrm{full}}+a\,A_{\mathrm{learned}},\qquad
A_{\mathrm{full}}=\mathbf 1\mathbf 1^\top-I,
\label{eq:blend}
\end{equation}
where $A_{\mathrm{learned}}$ is the strictly lower-triangular, binarized matrix of
\cref{app:code}. Three regimes follow.
\begin{itemize}[leftmargin=1.5em]
\item $a=0$ (\texttt{order\_mode} \texttt{"full"}): every latent is a possible parent of
      every other latent, with no order and no learned structure. This is not a DAG.
      The content prior is then a product of full conditionals, that is, a
      pseudo-likelihood and not a normalized density, so the KL values in this mode
      are not comparable with the ones after the switch.
\item $0<a<1$ (\texttt{"blend"}): the mixture of \cref{eq:blend}, which removes the
      jump between the two regimes; the order is introduced gradually and no reset of
      Adam is needed by default.
\item $a=1$ (\texttt{"tril"}): the lower-triangular DAG of the reference code. Only
      from here on do the sparsity penalty, the freezing and the sink-check gate
      have anything to act on; before it, the gate and \texttt{find\_sinks\_and\_fix}
      are inert.
\end{itemize}
With scheduling enabled the model starts at $a=0$; with scheduling disabled it starts
at $a=1$. In the $a=0$ regime the learnable adjacency receives no gradient, so the
probe of \cref{sec:schedule:gradlog} shows zero gradient for the adjacency terms.
This is expected.

\paragraph{When the blend starts, and when the next loss may start.}
The blend is part of the same chain of readiness as the losses. In order:
\begin{enumerate}[label=(\arabic*),leftmargin=2em]
\item REC reaches its plateau (the head reconstructs).
\item The KL ramp starts and reaches its target; KL is then \emph{ready}.
\item The blend ramp starts, at the end of the epoch in which KL becomes ready, and
      moves $a$ from $0$ to $1$ linearly over \texttt{schedule.tril\_ramp\_epochs}
      epochs ($a=0$ in the epoch after the start, as for the loss ramps; 
      $\texttt{schedule.tril\_ramp\_epochs} = 0$ means an immediate switch).
\item At $a=1$ the model calls the end-of-blend step, which resets the settle detector
      of KL.
\item \label{item:sche:sparsestart} SPARSE starts when KL has \emph{settled} after $a=1$. The wait is capped at
      \texttt{schedule.tril\_cap} epochs, counted from the epoch at which $a$ reached
      $1$; at the cap SPARSE starts anyway, with a warning, and the run becomes
      suspect.
\end{enumerate}
The reason for step \labelcref{item:sche:sparsestart} is that the switch from the full to the lower-
triangular 
adjacency changes the KL estimator, so KL jumps and then relaxes. SPARSE and the gate
must not see that transient: a gate that opens on a block that has not yet adapted
would call it stable for the wrong reason (cf.\ the opening rule of the gate below).
Consequently SPARSE starts at least \texttt{tril\_ramp\_epochs} plus
\texttt{settle\_window} epochs after KL is ready.

\paragraph{The settle detector.}
It is owned by the KL loss (\texttt{KLLoss.settled()}, \texttt{reset\_settle()}). It
observes the epoch mean of the validation KL when validation exists, otherwise of the
training KL, and fires when the relative change between \texttt{loss.settle\_window}
epochs ago and now is below \texttt{loss.settle\_tol}. It latches, and it is saved in
the checkpoint. The tolerance is a guess; a relative change is strict for KL values
near zero.

\paragraph{What this does not do.}
The blend does not choose the order. It postpones imposing it until the latents carry
information, and it makes the imposition smooth. A wrong order, and hence a false
sink, remains possible; the offline checks for it are listed in
\cref{sec:schedule:open}.

\begin{center}
\begin{tabular}{@{}lll@{}}
\toprule
phase & what is active & ends when\\
\midrule
A & REC, KL at its initial fraction; $a=0$ & REC plateau\\
B & KL ramp to target; $a=0$ & KL ready\\
C & blend ramp, $a:0\to1$ & $a=1$\\
D & lower-triangular DAG, SPARSE off & KL settled, or \texttt{tril\_cap}\\
E & SPARSE ramp, gate opens at \texttt{gate\_open\_frac} & all nodes frozen\\
F & fine-tune, graph fixed & \texttt{finetune\_epochs}\\
\bottomrule
\end{tabular}
\end{center}
BILIP starts at the REC plateau with its own cap, independently of this chain.

\subsection{Weight mechanics shared by all losses}
\label{sec:schedule:weights}
With scheduling enabled, the weight of a loss with target $\lambda$, initial fraction
$f_0$ ($0$ if none) and ramp length $R$ is
\[
w=\begin{cases}
0 & \text{switched off},\\
\lambda f_0 & \text{ramp not started},\\
\lambda\bigl(f_0+(1-f_0)\min(1,p/R)\bigr) & \text{ramp started, } p=\text{epochs
since start},
\end{cases}
\]
and $w=\lambda$ if $R\le0$. REC and INJ are created already started. A term whose
weight is $0$ at the moment is not computed (it is reported as $0$), so the expensive
BLAE terms cost nothing while inactive. The total is $L=\sum_k w_k\,\ell_k$ over REC,
KL, SPARSE, MORAL, INJ, BILIP, each $\ell_k$ unweighted. The order blend is not a
weight and is not part of this sum; it is a state of the prior.

\subsection{The six losses and the gate}
\label{sec:schedule:objects}

\paragraph{1. REC (constant, owns the plateau test).}
Weight $\lambda_{\mathrm{rec}}=10$ from epoch 1, never ramped or switched off (from
the toy setting of \cite{ngimpl}). REC is what the others wait for. Its plateau
detector fires once the epoch-mean REC has changed by less than $1\%$ (relative)
between $5$ epochs ago and now; the plateau latches. REC also answers whether the head
is learning at all: last epoch-mean REC below \texttt{rec\_frac}$=0.5$ times the
baseline. The baseline is the trivial reconstruction error, measured once by
\texttt{main.py} on the training rows as the variance of $x$ around its per-feature
mean, averaged over all entries (the MSE of the best constant predictor).

\paragraph{2. KL.}
\begin{itemize}[leftmargin=1.5em]
\item \emph{Stage A.} Weight $\texttt{kl\_init\_frac}\cdot\lambda_{\mathrm{kl}}^{\ast}$
      from epoch 1 ($10^{-3}$, a guess), so the head learns to reconstruct before the
      prior is allowed to pull on it.
\item \emph{Start.} The linear ramp (\texttt{kl\_ramp\_epochs}$=10$) to the target
      $\lambda_{\mathrm{kl}}^{\ast}=1$ starts when REC has plateaued. The estimator is
      the single-sample $\log q-\log p$ summed over the $D$ latents, since no analytic
      KL exists when the content prior mean depends on $z$.
\item \emph{Cap.} \texttt{kl\_cap}$=30$ epochs of waiting. At the cap: if REC is below
      $50\%$ of the baseline, start KL with a warning; if the baseline is unknown,
      start with a warning saying so; otherwise stop the run (the head is not
      learning).
\item \emph{Two readiness signals.} \texttt{ready()} is true when the ramp has reached
      the target; it starts the order blend. \texttt{settled()} is true when the
      settle detector of \cref{sec:schedule:phase0} has fired after $a=1$; it starts
      SPARSE.
\end{itemize}

\paragraph{3. SPARSE.}
Starts when KL has settled after the order blend has reached $1$
(\cref{sec:schedule:phase0}), with a $10$-epoch linear ramp from $0$ to a target of
$0.01$ (the toy value of \cite{ngimpl}). The penalty is the sum of
\cref{app:code:lossterms}, $\mathrm{sum}\bigl((I+A)^\top(I+A)\bigr)$ on the unresolved
block, divided by the fixed constant $n_c(n_c-1)/2$. The identity part is not
subtracted, so the value is never $0$ while the block is non-empty. It is switched off
when every content node is frozen. In disabled mode it has its target weight from
epoch 1, and $a=1$ from epoch 1.

\paragraph{4. The sink-check gate.}
The fixed \texttt{check\_epoch} clock of \cref{app:code} is replaced by a stability
gate.
\begin{itemize}[leftmargin=1.5em]
\item \emph{Inert before $a=1$.} In the full and blend modes the gate and the freezing
      do nothing, since there is no DAG yet to freeze.
\item \emph{Opening.} The gate is held closed until SPARSE has reached
      \texttt{gate\_open\_frac}$=0.5$ of its target weight, and then stays open.
      Before that the adjacency (initialized to all ones) has no pruning pressure and
      looks falsely stable. The gate's own counters start from the epoch it opens.
      This is a quick hack: the right fraction depends on the slope and the weight of
      SPARSE.
\item \emph{Stability.} After each epoch, compare the binarized unresolved block
      (entries above $0.5$, which is what the forward pass sees) with the previous
      epoch's. \texttt{stable\_run} counts consecutive epochs without change; it is
      reset on any flip and whenever the set of unfrozen nodes changes.
\item \emph{Freeze.} Denote $m$ the number of unfrozen nodes. When
      $\texttt{stable\_run}\ge W(m)$, freeze \emph{all} zero columns of the block at
      once (the nodes without remaining children), copy their rows and columns into
      \texttt{sure\_adj}, mask them in \texttt{sure\_mask}, log the epoch and node,
      and reset the counters. Starting $m$-dependent values:
\begin{center}
\begin{tabular}{@{}ccc@{}}
\toprule
$m$ (unfrozen) & $W$ & $\mathrm{cap}(m)$\\
\midrule
1--3 & 3 & 30\\
4--5 & 5 & 50\\
6 or more & 8 & 80\\
\bottomrule
\end{tabular}
\end{center}
\item \emph{Cap.} The cap counts epochs since the last reset. At the cap the run does
      \emph{not} freeze anything: it records the flip history, the raw entries nearest
      $0.5$ and the nodes frozen so far, and stops. No intervention (for example a
      hysteresis margin in the binarization) is added at first.
\item The strictly lower-triangular design forces the last-index node to have no
      children, so it is always frozen at the first successful check.
\item \emph{Deviation from the reference code.} \texttt{find\_sinknodes\_and\_fix} of
      \cref{app:code} indexes the global mask with indices taken from the compacted
      remaining block, which is correct only in the first round. Here compact indices
      are mapped back to global node ids, and the quirk that resets
      \texttt{should\_stop} when the adjacency sums to zero is dropped. The heuristic
      itself is unchanged and is still a mask heuristic, not the derivative test of
      the theory.
\item \emph{Clock mode} (scheduling off): freeze all zero columns every
      \texttt{check\_epoch} epochs, with no stability rule and no cap.
\end{itemize}

\paragraph{5. MORAL.}
Target $0$ for the first runs, in which case it is never computed. Once a target is
set, its ramp (10 epochs, a placeholder) starts at the first successful sink check.
After every check that froze $f$ nodes the model snapshots the binary adjacency as the
reference $A^{t-1}$ and sets $k_{\min}=n_c-f+1$; the loss sums the squared Frobenius
distances of $(I+A_{[:k,:k]})^\top(I+A_{[:k,:k]})$ to the same quantity of the
reference for $k=k_{\min},\dots,n_c$. Before the first snapshot it is $0$. Frozen
entries are constants and cancel in the difference, so the term measures the drift of
the unresolved block only, but it weights low-index unfrozen nodes more than the
original definition does (reweighting is deferred). The reference \texttt{prev\_adj}
and \texttt{k\_min} are plain attributes, not buffers, and are saved explicitly in the
schedule state. MORAL is switched off when all nodes are frozen.

\paragraph{6. INJ.}
The BLAE injective loss, computed on $\hat\mu$ and never on a sampled $\hat z$.
Constant weight $0.1$ from epoch 1 (a guess, no published anchor), threshold $0.3$,
never switched off. It needs a $(b,b)$ sub-table of the geodesic table for the
minibatch, so the trainer supplies one whenever the weight is positive.

\paragraph{7. BILIP.}
The BLAE decoder-Jacobian term, computed on $\hat\mu$, by far the most expensive loss
(about $D$ decoder passes per sample on a random fraction of the batch, with the plain
attention backend forced because the fused kernels have no forward-mode rule). It
starts when REC has plateaued, with a cap \texttt{bilip\_cap}$=60$ epochs, after which
it starts anyway with a warning. Linear ramp over $20$ epochs to a placeholder target
of $10^{-4}$. Its first computed value (one epoch after the ramp starts, the first
epoch with positive weight) is kept, so the target can be set against the real
magnitude. Upper bound $L=2$. The subset fraction is $0.3$ in the code against $10\%$
in the BLAE text (unresolved).

\subsection{The end-of-epoch sequence}
\label{sec:schedule:sequence}
The trainer calls \texttt{model.end\_epoch(epoch, stats)} once per epoch, after the
train and validation passes. \texttt{stats} holds the epoch means of the unweighted
terms. The fixed order is:
\begin{enumerate}[label=(\arabic*),leftmargin=2em]
\item every loss advances its ramp and observes \texttt{stats} (KL also feeds its
      settle detector);
\item the gate acts (subject to the opening rule above, and inert before $a=1$); if it
      froze nodes, MORAL takes its snapshot; a gate cap becomes a stop request;
\item only if scheduling is enabled, the rules, in this order: KL starts (REC plateau,
      cap 30); BILIP starts (REC plateau, cap 60); the order blend starts (KL ready);
      the order blend advances one step, and at $a=1$ the settle detector is reset;
      SPARSE starts (KL settled after $a=1$, cap \texttt{tril\_cap}); MORAL starts
      (first sink check); SPARSE and MORAL switch off (all nodes frozen);
\item run control decides whether to stop;
\item a \texttt{Status} is returned: stop flag and reason, suspect flag, this epoch's
      decisions and warnings, and the weights for the next epoch.
\end{enumerate}
Every decision and warning is appended to a log with its epoch, kept in the schedule
state, so a finished run shows why each loss started or stopped when it did. For a
healthy run the log shows, in this order: the KL start, the blend start, the end of
the blend, the SPARSE start, the opening of the gate, the freezes.

\subsection{Run control, suspect runs and outcomes}
\label{sec:schedule:control}
\begin{itemize}[leftmargin=1.5em]
\item \texttt{max\_epochs}$=250$ is a hard stop. The earlier estimate of a worst-case
      budget of about $190$ epochs at $n_c=5$ belonged to an earlier design and has
      not been recomputed; the blend and the settle wait add at least
      \texttt{tril\_ramp\_epochs} plus \texttt{settle\_window} epochs.
\item When all content nodes are frozen, continue for \texttt{finetune\_epochs} (a
      guess of $20$) with the graph fixed and SPARSE and MORAL off, then stop with
      ``finished''.
\item Stop priority: a stop request (KL cap with the head not learning, or sink-check
      cap), then the end of the fine-tune phase (scheduling enabled only), then
      \texttt{max\_epochs}.
\item A run is \emph{suspect} if any cap was hit (including \texttt{tril\_cap}), or
      (scheduling enabled) it stopped before all nodes froze. The first epoch at which
      this becomes true is logged as \texttt{SUSPECT}; the run continues but must not
      be trusted.
\item \textbf{Outcomes} written to \texttt{report.json}: \texttt{finished},
      \texttt{finished\_suspect}, \texttt{max\_epochs} (scheduling disabled), and
      \texttt{failure:<kind>} with kind in \texttt{nonfinite\_loss},
      \texttt{model\_stop}, or \linebreak
      \texttt{max\_epochs\_unfinished} (scheduling enabled only). 
      On a failure the trainer also writes \linebreak
      \texttt{failure.json} with the
      reason, epoch, caps hit, freeze log, gate diagnostics, last stats, current
      weights, the decision and warning logs and the model report. \texttt{main.py}
      exits with status $1$ on a failure outcome.
\item The report also contains the final \texttt{order\_mode} and \texttt{order\_blend}
      and the epochs at which the blend started and reached $1$. These, together with
      the checkpointed blend position, make a resumed run continue the blend where it
      stopped.
\end{itemize}

\subsection{Gradient-scale probe}
\label{sec:schedule:gradlog}
The weights are tuned with the Adam-based heuristic
\begin{equation}
\mathrm{scaler}=
\frac{0.01}{\mathrm{LR}\times\texttt{weight}\times\lVert\nabla_{\texttt{unweighted}}\rVert}=
\frac{0.01}{\mathrm{LR}\times\lVert\nabla_{\texttt{weighted}}\rVert},
\label{eq:scaler}
\end{equation}
tuning each weight until its scaler is about $1$. The probe (\texttt{probe.py}) is
diagnostic only: it never changes a weight, never stops a run and never touches an
optimizer.
\begin{itemize}[leftmargin=1.5em]
\item \textbf{Norm.} For a term $k$ with unweighted loss $\ell_k$ and declared
      parameter set $P_k$,
      \[
      \mathrm{norm}_k=\Bigl(\sum_{p\in P_k}\mathrm{LR}_p^2\,\bigl\lVert\partial\ell_k
      /\partial p\bigr\rVert^2\Bigr)^{1/2},
      \]
      with $\mathrm{LR}_p$ the learning rate of the optimizer group that owns $p$. The
      weighted norm is \linebreak 
      $w_k \cdot \mathrm{norm}_k$, linear in the weight, so
      $\mathrm{scaler}_k(w)=0.01/(w \cdot \mathrm{norm}_k)$ and the suggested weight is
      $w^{\ast}_k=0.01/\mathrm{norm}_k$ (for which $\mathrm{scaler}_k(w)=1$).
\item \textbf{Computed at the target weight.} A ramping term is smaller by design, so
      the scaler is reported at the target weight, with the ramp fraction shown next
      to it.
\item \textbf{Flags only.} ``strong'' if the scaler is below $0.1$, ``weak'' if above
      $10$ (a loose band), plus ``no\_gradient'' and ``target\_zero''. A flag is
      logged as \texttt{PROBE-FLAG}; there is no automatic stop, including for
      blow-ups. While $a=0$ the adjacency terms show ``no\_gradient'' by construction
      (\cref{sec:schedule:phase0}).
\item \textbf{Declared parameter sets} (\texttt{PROBE\_SETS}, proposal, to be
      confirmed):
\begin{center}
\begin{tabular}{@{}ll@{}}
\toprule
term & parameters\\
\midrule
REC & whole head (encoder side and decoder side)\\
KL & encoder side, prior, adjacency\\
SPARSE, MORAL & adjacency\\
INJ & encoder side\\
BILIP & decoder side and encoder side\\
\bottomrule
\end{tabular}
\end{center}
      Encoder side means the \texttt{AutoencoderKL} encoder, \texttt{quant\_conv} and
      \texttt{pre\_linear}; decoder side means its decoder,
      \texttt{post\_quant\_conv} and \texttt{post\_linear}. A parameter that fits no
      group raises an error, so a renamed module is caught rather than silently
      dropped.
\item \textbf{Computation.} Every \texttt{probe.every}$=5$ epochs (BILIP only every
      \texttt{probe.bilip\_every}$=25$), after \texttt{end\_epoch} and before the
      checkpoint, on \texttt{probe.n\_batches} training batches. One forward pass per
      batch, then \texttt{torch.autograd.grad} per term with
      \texttt{retain\_graph=True}, which does not write \texttt{.grad}, so a training
      step cannot be disturbed. Terms with weight $0$ are included. The RNG state, the
      train/eval mode and BILIP's recorded first value are restored afterwards. A
      probe that raises is logged as a warning and training continues.
\item \textbf{Skipped terms}: SPARSE with no unresolved block, MORAL before the first
      sink check, INJ without a geodesic table, BILIP when not due.
\item \textbf{Adjacency.} Its optimizer is SGD with momentum, where the step is really
      LR times the gradient, whereas \cref{eq:scaler} was calibrated with Adam. The
      SPARSE and MORAL rows are therefore marked unvalidated.
\end{itemize}
Output: a table on stdout, one JSON line per probe in \texttt{gradscale.jsonl}, and
the flag lines in \texttt{events.log}.

\subsection{What the trainer does}
\label{sec:schedule:trainer}
\begin{itemize}[leftmargin=1.5em]
\item \textbf{Two optimizers, disjoint.} Adam on all non-adjacency parameters
      (\texttt{lr}$=10^{-3}$), SGD with momentum $0.9$ on the adjacency
      (\texttt{adj\_lr}$=10^{-3}$). This fixes the overlap of \cref{app:code}, where
      the adjacency was stepped by both. Learning rates are constant; LR scheduling
      is parked, and the adjacency LR in particular should stay constant. No optimizer
      reset is made at the phase transitions: the blend removes the jump that a reset
      would address.
\item \textbf{One epoch.} Train pass (one optimizer step per batch, epoch means of
      every unweighted term), validation pass, \texttt{model.end\_epoch}, logging of
      decisions and warnings, probe, checkpoint, stop if requested.
\item \textbf{Validation.} \texttt{val\_rec} and \texttt{val\_kl} with $z$ sampled as
      in training, and, optionally, macro species accuracy: species predicted from the
      content latents only, with the time of day observed, as
      $\arg\max_y\log p(\hat\mu_c\mid y,t,e)$ under the current prior with a uniform
      label prior, $z=\hat\mu$. The accuracy costs one prior pass per species per
      batch and is switched off by \texttt{--no\_val\_acc}. It is an addition to the
      plan, kept pending confirmation. \texttt{val\_kl} also feeds the settle detector
      of \cref{sec:schedule:phase0}.
\item \textbf{Non-finite loss.} A NaN or infinite term raises before \texttt{backward},
      so no optimizer step happens, the run stops with outcome
      \texttt{failure:nonfinite\_loss}, and the last good checkpoint is not
      overwritten.
\item \textbf{Checkpoint, at epoch boundaries only}, written atomically as
      \texttt{last.pt}: model state (including the buffers \texttt{sure\_mask} and
      \texttt{sure\_adj}), the schedule state (\texttt{stable\_run}, epochs since the
      last reset, the previous epoch's binarized block, \texttt{prev\_adj},
      \texttt{k\_min}, the REC window, the KL settle detector, the order blend and its
      ramp position, stage and ramp positions, run control, the decision logs), both
      optimizer states including the SGD momentum, all RNG states, the history, the
      finalized \texttt{RunConfig}, the trainer configuration and the input
      arguments. \texttt{--resume} restores everything and continues from the next
      epoch, using the saved arguments and \texttt{RunConfig}; only \texttt{--device}
      and new \texttt{--set} overrides come from the command line. Older checkpoints
      without the blend state still load.
\item \textbf{Geodesic table.} Built once, only if the injective loss is active, on all
      rows of the selected training locations (train and validation rows together, so
      the validation rows act as extra geometry for the regularizer targets).
\end{itemize}

\subsection{Configuration added for phase 0}
\label{sec:schedule:newcfg}
\begin{center}
\begin{tabular}{@{}lcl@{}}
\toprule
field & default & meaning\\
\midrule
\texttt{loss.settle\_window} & 5 & epochs compared by the KL settle detector\\
\texttt{loss.settle\_tol} & 0.02 & relative change below which KL has settled\\
\texttt{schedule.tril\_ramp\_epochs} & 10 & length of the blend ramp (0: immediate)\\
\texttt{schedule.tril\_cap} & 60 & max wait for the settle after $a=1$\\
\bottomrule
\end{tabular}
\end{center}
All four are validated in \texttt{config.py}, and all four are guesses.

\subsection{Open items}
\label{sec:schedule:open}
\begin{itemize}[leftmargin=1.5em]
\item \textbf{To confirm:} \texttt{PROBE\_SETS}; keeping the validation accuracy; the
      per-feature versus the global variance for the REC baseline; building the
      geodesic table on training rows only (which would need the ids remapped);
      \texttt{settle\_tol}.
\item \textbf{Phase 0:} the KL values while $a<1$ come from a pseudo-likelihood and are
      not comparable with later ones. The first real runs should show whether the
      settle detector fires within \texttt{tril\_cap}.
\item \textbf{Offline checks of the final graph} (no influence on the schedule): macro
      accuracy; a sanity check of the graph for degenerate structure; embeddings of
      $\hat\mu$ (t-SNE or UMAP, and a direct 2D scatter of the style latents), read
      only as separated versus mixed; leave-one-location-out over the training
      locations to see style leaking into content (L100 stays untouched); comparison
      of graphs across seeds, matching latents by the correlation of $\hat\mu$.
\item \textbf{Not started:} LR scheduling (none or cosine, applied by the trainer,
      adjacency LR constant by default); an optimizer reset for the prior parameters
      at a transition, as a separate switch, only if the probe shows a jump;
      post-hoc order selection by scoring permutations of the content slots; a
      comparison with \texttt{kl\_init\_frac}$=0$ across seeds. \texttt{TrainerConfig}
      still lives in \texttt{trainer.py} with its own defaults, against the rule that
      defaults live only in \texttt{config.py}, and has to move.
\item \textbf{Gate:} a margin in the binarization and any intervention at the
      sink-check cap; \texttt{gate\_open\_frac} may be wrong for slow ramps or small
      weights.
\item \textbf{Deferred:} normalization of KL, MORAL and BILIP across $n_c$; MORAL
      reweighting; validation split by \texttt{src\_uuid} when synthetic images are
      included; the BILIP subset fraction.
\item \textbf{Untested:} the torch-dependent code, in this order: the smoke test
      (\texttt{python main.py --synthetic --set schedule.max\_epochs=4 --set
      loss.lambda\_inj=0}), \texttt{python m1net.py} (blend $0$, $0.5$, $1$), the
      gate's tensor code, the loss computations, \texttt{M1Model.\_\_init\_\_}, and
      one probe (\texttt{--set probe.every=1}). For a longer smoke run use small
      values, for example \texttt{--set schedule.tril\_ramp\_epochs=2 --set
      loss.settle\_window=2}, and check that the decision log shows the blend start,
      the blend end and the SPARSE start in that order.
\item Every starting value above is a guess; all are to be fixed from the first runs.
\end{itemize}
